\documentclass[a4paper,12pt]{article}
\usepackage[T2A]{fontenc}
\usepackage[utf8]{inputenc}
\usepackage[russian]{babel}
\usepackage{amsmath,amssymb,mathtools}
\renewcommand{\sfdefault}{cmss}
\renewcommand{\familydefault}{\sfdefault}
\usepackage{icomma}
\usepackage[margin=23mm,headheight=23pt,footskip=16mm]{geometry}
\usepackage{microtype}
\usepackage{tikz}
\usetikzlibrary{calc,arrows.meta,positioning,shapes.geometric}
\usepackage{enumitem}
\usepackage{tabularx,booktabs}
\usepackage{fancyhdr}
\usepackage{setspace}
\usepackage{xcolor}
\usepackage{hyperref}

\definecolor{accent}{HTML}{087B7B}
\definecolor{darkaccent}{HTML}{075455}
\definecolor{textgray}{HTML}{343A3D}
\definecolor{softgray}{HTML}{707D80}
\definecolor{fine}{HTML}{D5DFDF}
\hypersetup{colorlinks=true,linkcolor=darkaccent,pdfauthor={Лёвин Артём Александрович},
pdftitle={Чек-лист. Метод Крамера и метод Гаусса}}
\setstretch{1.18}
\setlength{\parindent}{0pt}
\setlength{\parskip}{0.53em}
\setlist[enumerate]{leftmargin=*,itemsep=0.65em,topsep=0.5em}
\pagestyle{fancy}
\fancyhf{}
\fancyhead[L]{\footnotesize\color{softgray} Математика \,\textbullet\, 10 класс}
\fancyhead[R]{\footnotesize\color{softgray} 18.10.26}
\fancyfoot[C]{\footnotesize\color{softgray}\thepage}
\renewcommand{\headrulewidth}{0.35pt}
\renewcommand{\headrule}{\hbox to\headwidth{\color{fine}\leaders\hrule height \headrulewidth\hfill}}
\newcommand{\subnote}[1]{{\color{softgray}\small #1}}
\newcommand{\strong}[1]{{\color{darkaccent}\textbf{#1}}}
\newcommand{\task}[1]{\item #1}
\newcommand{\aug}[2]{\left(\begin{array}{#1|r}#2\end{array}\right)}
\tikzset{
  flowstart/.style={ellipse,draw=darkaccent,very thick,align=center,
    minimum width=32mm,minimum height=10mm,fill=white,font=\small},
  flowact/.style={rounded corners=2pt,rectangle,draw=accent,very thick,
    align=center,text width=82mm,minimum height=15mm,inner sep=6pt,font=\small},
  flowwide/.style={rounded corners=2pt,rectangle,draw=accent,very thick,
    align=center,text width=100mm,minimum height=15mm,inner sep=6pt,font=\small},
  flowcond/.style={diamond,aspect=2.4,draw=darkaccent,very thick,
    align=center,inner sep=2pt,text width=34mm,font=\small},
  arr/.style={-{Latex[length=3mm,width=2mm]},thick,color=darkaccent}
}

\begin{document}
\thispagestyle{empty}
\vspace*{32mm}
\begin{center}
{\color{accent}\large МАТЕМАТИКА \,\textbullet\, 10 КЛАСС}\\[14mm]
{\color{softgray}\large Персональный учебный материал}\\[8mm]
{\Huge\bfseries Чек-лист}\\[8mm]
{\LARGE\bfseries\color{darkaccent} Метод Крамера\\[2mm] и метод Гаусса}\\[15mm]
{\Large Системы линейных уравнений}\\[19mm]
{\color{textgray}\large 18 октября 2026 года}
\end{center}
\vfill
\begin{center}
{\color{darkaccent}\large Лёвин Артём Александрович}\\[2mm]
{\large эксклюзивно для Анны Трапезниковой}
\end{center}
\vspace{14mm}
\begin{tikzpicture}[remember picture,overlay]
\draw[accent!60,line width=1.0pt]
([xshift=23mm,yshift=19mm]current page.south west)
.. controls ([xshift=65mm,yshift=32mm]current page.south west)
and ([xshift=-64mm,yshift=5mm]current page.south east)
.. ([xshift=-23mm,yshift=20mm]current page.south east);
\end{tikzpicture}
\clearpage

\section*{\color{darkaccent}01.\quad Что необходимо понимать}
Систему двух линейных уравнений можно записать как
\[
\begin{cases}
a_1x+b_1y=c_1,\\
a_2x+b_2y=c_2.
\end{cases}
\]
\strong{Совместная система} имеет хотя бы одно решение. Если решений нет, система \strong{несовместна}.

\textbf{Проверочный список}
\begin{enumerate}[label=$\square$,leftmargin=11mm,itemsep=0.3em]
\task Различаю основную матрицу коэффициентов и столбец свободных членов.
\task Умею вычислить определитель второго порядка:
\[
\Delta=\begin{vmatrix}a&b\\c&d\end{vmatrix}=ad-bc.
\]
\task Помню: равенство нулю главного определителя требует \emph{дополнительного анализа}.
\task Понимаю смысл элементарных преобразований строк: они сохраняют множество решений.
\task Умею восстановить уравнения по строкам расширенной матрицы.
\end{enumerate}

\section*{\color{darkaccent}02.\quad Метод Крамера}
Для системы двух уравнений вычислим три определителя:
\[
\Delta=\begin{vmatrix}a_1&b_1\\a_2&b_2\end{vmatrix},\qquad
\Delta_x=\begin{vmatrix}c_1&b_1\\c_2&b_2\end{vmatrix},\qquad
\Delta_y=\begin{vmatrix}a_1&c_1\\a_2&c_2\end{vmatrix}.
\]
В $\Delta_x$ заменяем \textbf{столбец при $x$}, в $\Delta_y$ ---
\textbf{столбец при $y$}; порядок строк сохраняем.

\subnote{Нельзя делить на $\Delta$, если $\Delta=0$.
Нулевые частные определители при $\Delta\ne0$ допустимы.}

\begin{center}
\begin{tabularx}{\linewidth}{@{}p{46mm}X@{}}
\toprule
Условие & Вывод\\
\midrule
$\Delta\ne0$ & Единственное решение:
$x=\dfrac{\Delta_x}{\Delta}$,\enspace
$y=\dfrac{\Delta_y}{\Delta}$.\\[2mm]
$\Delta=0$; хотя бы один из $\Delta_x,\Delta_y$ ненулевой
& Решений нет.\\
$\Delta=\Delta_x=\Delta_y=0$ & Бесконечно много решений
(для системы из двух уравнений с двумя неизвестными).\\
\bottomrule
\end{tabularx}
\end{center}

\clearpage
\section*{\color{darkaccent}03.\quad Два показательных случая}
\textbf{Единственное решение:}
\[
\begin{cases}2x+y=5,\\x-y=1.\end{cases}
\quad
\Delta=-3,\quad\Delta_x=-6,\quad\Delta_y=-3.
\quad\Rightarrow\quad (x,y)=(2,1).
\]

\textbf{Решений нет:}
\[
\begin{cases}-7x-y=13,\\21x+3y=0.\end{cases}
\quad
\Delta=0,\quad\Delta_x=39,\quad\Delta_y=-273.
\]
Поскольку $\Delta=0$, но $\Delta_x\ne0$, система несовместна.

\strong{Однородная система}: справа во всех уравнениях стоят нули.
Решение, в котором все неизвестные равны нулю, существует всегда.
Если $\Delta\ne0$, оно единственное; если $\Delta=0$, есть также
ненулевые решения.

\clearpage
\section*{\color{darkaccent}Алгоритм 1.\quad Метод Крамера}
\vspace{8mm}
\begin{center}
\begin{tikzpicture}[node distance=11mm and 22mm]
\node[flowstart] (a) {Начало};
\node[flowact,below=of a] (b) {Вычислить $\Delta$, $\Delta_x$, $\Delta_y$};
\node[flowcond,below=13mm of b] (c) {$\Delta\ne0$?};
\node[flowact,text width=38mm] (d) at ($(c.south)+(-4.1cm,-3.0cm)$)
{$x=\Delta_x/\Delta$\\$y=\Delta_y/\Delta$};
\node[flowcond] (e) at ($(c.south)+(2.6cm,-3.0cm)$)
{$\Delta_x=0$\\ и $\Delta_y=0$?};
\node[flowact,text width=35mm] (f) at ($(e.south)+(-2.3cm,-2.6cm)$)
{Бесконечно\\ много решений};
\node[flowact,text width=34mm] (g) at ($(e.south)+(2.3cm,-2.6cm)$)
{Решений нет};
\draw[arr] (a)--(b);
\draw[arr] (b)--(c);
\draw[arr] (c.west) -| node[pos=0.26,above,font=\small]{да} (d.north);
\draw[arr] (c.south) |- node[pos=0.26,right,font=\small]{нет} (e.west);
\draw[arr] (e.south) -| node[pos=0.25,above,font=\small]{да} (f.north);
\draw[arr] (e.south) -| node[pos=0.25,above,font=\small]{нет} (g.north);
\end{tikzpicture}
\end{center}
\clearpage

\section*{\color{darkaccent}04.\quad Метод Гаусса: главная идея}
Систему из трёх уравнений записывают \strong{расширенной матрицей}:
после вертикальной черты помещают свободные члены.
\[
\left(
\begin{array}{rrr|r}
2&3&-5&4\\
3&-4&2&21\\
7&8&-6&1
\end{array}
\right).
\]
Нужно последовательно получить \strong{ступенчатый вид}:
нули под ведущими элементами строк. Затем найти неизвестные
\strong{обратным ходом} --- от последней строки к первой.

\textbf{Допустимые преобразования строк}
\begin{enumerate}[label=\arabic*.,itemsep=0.25em]
\task Переставить две строки.
\task Умножить строку на ненулевое число.
\task Прибавить к одной строке другую, умноженную на любое число.
\end{enumerate}

\textbf{Разбор того же примера.} Сначала обнуляем элементы под
первой двойкой. Обозначение $R_i$ означает строку с номером $i$.
\[
\begin{aligned}
R_2&\leftarrow 2R_2-3R_1,\\
R_3&\leftarrow 2R_3-7R_1.
\end{aligned}
\qquad\Rightarrow\qquad
\left(
\begin{array}{rrr|r}
2&3&-5&4\\
0&-17&19&30\\
0&-5&23&-26
\end{array}
\right).
\]
В частности, второй новый ряд равен
\[
(6,-8,4\mid42)-(6,9,-15\mid12)
=(0,-17,19\mid30).
\]
Для обнуления $-5$ в третьей строке приводим $-17$ и $-5$
к общему значению $-85$:
\[
R_3\leftarrow 17R_3-5R_2.
\]
\[
\Rightarrow\qquad
\left(
\begin{array}{rrr|r}
2&3&-5&4\\
0&-17&19&30\\
0&0&296&-592
\end{array}
\right).
\]
\textbf{Обратный ход:}
\[
\begin{aligned}
296z&=-592 &&\Rightarrow z=-2,\\
-17y+19(-2)&=30 &&\Rightarrow y=-4,\\
2x+3(-4)-5(-2)&=4 &&\Rightarrow x=3.
\end{aligned}
\]
\strong{Ответ}: $(x,y,z)=(3,-4,-2)$.
Проверка подстановкой в первоначальные три уравнения даёт
соответственно $4$, $21$, $1$.

\section*{\color{darkaccent}05.\quad На что обращать внимание}
\begin{enumerate}[label=$\square$,leftmargin=11mm,itemsep=0.35em]
\task Выполняю действие \emph{со всеми} элементами строки,
включая число справа от черты.
\task Заменяю именно ту строку, для которой вычисляю новую запись.
\task Выбираю множители так, чтобы нужные коэффициенты сократились.
\task В выражении $-12-(-35)$ получаю $23$, учитывая оба минуса.
\task Не считаю дробный ответ признаком ошибки.
\task Если появилась строка $(0,0,0\mid c)$ при $c\ne0$,
записываю противоречие $0=c$: решений нет.
\task Если появляется строка $(0,0,0\mid0)$,
проверяю количество свободных переменных: возможны бесконечно многие решения.
\end{enumerate}

\clearpage
\section*{\color{darkaccent}Алгоритм 2.\quad Прямой и обратный ход Гаусса}
\vspace{5mm}
\begin{center}
\begin{tikzpicture}[node distance=9mm]
\node[flowstart] (a) {Начало};
\node[flowwide,below=of a] (b) {Записать расширенную матрицу};
\node[flowwide,below=of b] (c)
{Выбрать ведущий ненулевой коэффициент.\\При необходимости переставить строки};
\node[flowwide,below=of c] (d)
{Обнулить элементы \emph{под} ведущим элементом,
выполняя преобразования целых строк};
\node[flowwide,below=of d] (e)
{Повторить для следующих столбцов,\\получить ступенчатый вид};
\node[flowcond,below=13mm of e] (f)
{Есть строка\\$0=c$, $c\ne0$?};
\node[flowact,right=20mm of f,text width=32mm] (g) {Решений нет};
\node[flowwide,below=19mm of f] (h)
{Обратный ход: найти ведущие переменные снизу вверх};
\node[flowwide,below=of h] (i)
{Если есть свободные переменные, задать параметры;
иначе получить единственное решение};
\node[flowstart,below=of i] (j) {Проверка};
\draw[arr] (a)--(b);
\draw[arr] (b)--(c);
\draw[arr] (c)--(d);
\draw[arr] (d)--(e);
\draw[arr] (e)--(f);
\draw[arr] (f) -- node[above,font=\small]{да} (g);
\draw[arr] (f) -- node[right,font=\small]{нет} (h);
\draw[arr] (h)--(i);
\draw[arr] (i)--(j);
\end{tikzpicture}
\end{center}
\clearpage

\section*{\color{darkaccent}06.\quad Самостоятельная тренировка}
\textbf{Путь А \,\textbullet\, 10 заданий}

\subnote{Решайте последовательно. Для заданий 1--5 используйте метод Крамера;
для заданий 6--10 --- метод Гаусса. При отсутствии единственного решения
обязательно укажите, сколько решений имеет система.}

\subsection*{\color{darkaccent}Часть I. Метод Крамера}
\renewcommand{\arraystretch}{1.65}
\begin{tabular}{@{}p{.48\linewidth}p{.48\linewidth}@{}}
\textbf{1.}\quad $\displaystyle\begin{cases}2x+y=5,\\x-y=1.\end{cases}$ &
\textbf{2.}\quad $\displaystyle\begin{cases}3x-2y=5,\\x+2y=7.\end{cases}$\\
\textbf{3.}\quad $\displaystyle\begin{cases}2x-y=0,\\3x+4y=0.\end{cases}$ &
\textbf{4.}\quad $\displaystyle\begin{cases}-7x-y=13,\\21x+3y=0.\end{cases}$\\
\textbf{5.}\quad $\displaystyle\begin{cases}2x+y=6,\\4x+2y=12.\end{cases}$ & \\
\end{tabular}
\renewcommand{\arraystretch}{1.0}

\subsection*{\color{darkaccent}Часть II. Метод Гаусса}
\begin{enumerate}[label=\textbf{\arabic*.},start=6,itemsep=0.4em]
\task $\displaystyle\begin{cases}
x+y+z=6,\\2x-y+z=3,\\-x+2y+3z=12.
\end{cases}$
\task $\displaystyle\begin{cases}
x+y-z=4,\\2x-y+3z=0,\\-x+3y-2z=3.
\end{cases}$
\task $\displaystyle\begin{cases}
2y+z=4,\\x-y+z=3,\\2x+y-z=3.
\end{cases}$
\task $\displaystyle\begin{cases}
x+y+z=0,\\2x+2y+2z=0,\\x-y+z=0.
\end{cases}$
\task $\displaystyle\begin{cases}
x+y+z=3,\\2x+2y+2z=7,\\x-y+z=1.
\end{cases}$
\end{enumerate}
\clearpage

\section*{\color{darkaccent}07.\quad Ответы для самопроверки}
\subnote{Сначала завершите вычисления самостоятельно; затем сверьте итог.
В системах с бесконечно многими решениями параметр $t$ --- произвольное
действительное число.}

\vspace{6mm}
\renewcommand{\arraystretch}{1.55}
\begin{tabularx}{\linewidth}{@{}p{15mm}X@{}}
\toprule
№ & Ответ\\
\midrule
1 & $(x,y)=(2,1)$.\\
2 & $(x,y)=(3,2)$.\\
3 & $(x,y)=(0,0)$; единственное решение.\\
4 & Решений нет.\\
5 & Бесконечно много решений:
$(x,y)=\left(3-\dfrac{t}{2},\,t\right)$, $t\in\mathbb R$.\\
6 & $(x,y,z)=(1,2,3)$.\\
7 & $(x,y,z)=(2,1,-1)$.\\
8 & $(x,y,z)=(2,1,2)$.\\
9 & Бесконечно много решений:
$(x,y,z)=(-t,0,t)$, $t\in\mathbb R$.\\
10 & Решений нет.\\
\bottomrule
\end{tabularx}

\vspace{10mm}
\textbf{Ключевые самопроверки}

\begin{itemize}[leftmargin=7mm,itemsep=0.55em]
\item В задании 3: $\Delta=11\ne0$ и
$\Delta_x=\Delta_y=0$. Поэтому $(0,0)$ --- \emph{единственное} решение.
\item В задании 4: $\Delta=0$, но $\Delta_x=39$.
\item В задании 5: $\Delta=\Delta_x=\Delta_y=0$.
\item В задании 9 последняя строка при исключении переменных
превращается в тождество $0=0$.
\item В задании 10 второе уравнение противоречит первому:
$2(x+y+z)=6$, а в условии стоит $7$.
\end{itemize}
\end{document}
